\documentclass[12pt,a4paper]{article}

% Packages
\usepackage[utf8]{inputenc}
\usepackage[T1]{fontenc}
\usepackage{amsmath,amssymb,amsfonts}
\usepackage{graphicx}
\usepackage{booktabs}
\usepackage{hyperref}
\usepackage{listings}
\usepackage{xcolor}
\usepackage{geometry}
\usepackage{titlesec}
\usepackage{enumitem}
\usepackage{fancyhdr}
\usepackage{caption}
\usepackage{subcaption}

% Page geometry
\geometry{margin=1in}

% Hyperlink setup
\hypersetup{
    colorlinks=true,
    linkcolor=blue,
    filecolor=magenta,
    urlcolor=cyan,
    citecolor=blue,
}

% Code listing style
\definecolor{codegreen}{rgb}{0,0.6,0}
\definecolor{codegray}{rgb}{0.5,0.5,0.5}
\definecolor{codepurple}{rgb}{0.58,0,0.82}
\definecolor{backcolour}{rgb}{0.95,0.95,0.92}

\lstdefinestyle{mystyle}{
    backgroundcolor=\color{backcolour},
    commentstyle=\color{codegreen},
    keywordstyle=\color{magenta},
    numberstyle=\tiny\color{codegray},
    stringstyle=\color{codepurple},
    basicstyle=\ttfamily\footnotesize,
    breakatwhitespace=false,
    breaklines=true,
    captionpos=b,
    keepspaces=true,
    numbers=left,
    numbersep=5pt,
    showspaces=false,
    showstringspaces=false,
    showtabs=false,
    tabsize=2,
    frame=single,
    framerule=0.5pt,
}

\lstset{style=mystyle}

% Section formatting
\titleformat{\section}{\Large\bfseries}{\thesection}{1em}{}
\titleformat{\subsection}{\large\bfseries}{\thesubsection}{1em}{}
\titleformat{\subsubsection}{\normalsize\bfseries}{\thesubsubsection}{1em}{}

% Header and footer
\pagestyle{fancy}
\fancyhf{}
\fancyhead[L]{SynBO: Bayesian Optimization for Chemical Reactions}
\fancyhead[R]{\thepage}

% Title information
\title{\textbf{\Large SynBO: Synthetic Bayesian Optimization for\\ Chemical Reaction Condition Screening}\\[1em]
\large An Intelligent Tool for Synthetic Chemists}
\author{Zhenzhi Tan\\[0.5em]
\texttt{zhenzhi-tan@outlook.com}}
\date{\today}

\begin{document}

\maketitle

\begin{abstract}
\noindent
\textbf{SynBO} (Synthetic Bayesian Optimization) is an intelligent reaction optimization tool designed specifically for synthetic chemists. It leverages Bayesian Optimization (BO) algorithms to help researchers find optimal reaction conditions with minimal experimental effort. This article provides a comprehensive introduction to SynBO, explaining its core principles, architecture, and practical workflow. We demonstrate how chemists can efficiently utilize SynBO through an integrated skill-based interface to accelerate reaction optimization, reduce experimental costs, and discover optimal conditions typically within 50--80 experiments instead of hundreds. The tool supports both single-objective (e.g., yield maximization) and multi-objective optimization (e.g., yield and enantioselectivity), making it applicable to a wide range of synthetic chemistry challenges.

\vspace{1em}
\noindent\textbf{Keywords:} Bayesian Optimization, Reaction Screening, Machine Learning, Synthetic Chemistry, High-Throughput Experimentation
\end{abstract}

\tableofcontents
\newpage

% ============================================================================
\section{Introduction}
% ============================================================================

\subsection{The Challenge of Reaction Optimization}

Optimizing a new chemical reaction is one of the most time-consuming and resource-intensive tasks in synthetic chemistry. A typical reaction optimization process involves screening numerous combinations of reaction conditions:

\begin{itemize}
    \item \textbf{Catalysts}: Various organocatalysts or metal complexes
    \item \textbf{Solvents}: Polar, non-polar, protic, aprotic options
    \item \textbf{Bases/Additives}: Acids, bases, ligands, electrolytes
    \item \textbf{Temperature}: From cryogenic to elevated temperatures
    \item \textbf{Concentration}: Reactant and reagent concentrations
    \item \textbf{Reaction Time}: From minutes to days
\end{itemize}

The traditional approach, known as \textbf{OFAT} (One-Factor-At-A-Time), is fundamentally inefficient. Consider a typical scenario with 5 catalysts, 5 solvents, 4 bases, and 4 temperatures: this yields $5 \times 5 \times 4 \times 4 = 400$ combinations. Testing all combinations is clearly impractical in terms of time, materials, and labor.

\subsection{The SynBO Solution}

\textbf{SynBO} addresses this challenge by employing Bayesian Optimization (BO), a machine learning technique that intelligently guides experimental design. Like an experienced chemist, SynBO ``learns'' from previous experiments and ``predicts'' which conditions are most likely to succeed next. The key advantages include:

\begin{enumerate}
    \item \textbf{Efficiency}: Typically requires only 50--80 experiments to find optimal conditions
    \item \textbf{Intelligence}: Balances exploitation (trying promising conditions) and exploration (testing unknown areas)
    \item \textbf{Flexibility}: Supports single and multi-objective optimization
    \item \textbf{Accessibility}: User-friendly interface designed for chemists without machine learning expertise
\end{enumerate}

\subsection{Article Organization}

This article is organized as follows: Section~\ref{sec:bo_basics} introduces the fundamentals of Bayesian Optimization with chemistry analogies. Section~\ref{sec:architecture} describes the software architecture and core components. Section~\ref{sec:workflow} presents the practical workflow for chemists using SynBO skills. Section~\ref{sec:algorithms} details the underlying algorithms and their chemical interpretations. Section~\ref{sec:case_studies} showcases real-world applications. Finally, Section~\ref{sec:conclusion} concludes with future perspectives.

% ============================================================================
\section{Bayesian Optimization Basics}
\label{sec:bo_basics}
% ============================================================================

\subsection{The Mountain Climbing Analogy}

Bayesian Optimization can be intuitively understood through a mountain climbing analogy. Imagine you are searching for the highest peak in complete darkness:

\begin{enumerate}
    \item \textbf{Initialization}: Take a few random steps and record the altitude. This corresponds to running a small number of initial experiments (typically 5--10) and recording yields or selectivities.
    
    \item \textbf{Build a Mental Map}: Based on where you've been, infer the shape of the entire mountain range. The algorithm learns reaction patterns from experimental data using surrogate models.
    
    \item \textbf{Intelligent Decision}: Decide where to go next---either toward areas that might be higher (\textit{exploitation}) or explore completely unknown areas (\textit{exploration}).
    
    \item \textbf{Iterate}: Repeat steps 2--3 until you find the highest peak, corresponding to finding optimal reaction conditions.
\end{enumerate}

\subsection{Chemistry Interpretation}

The BO process mirrors how experienced chemists optimize reactions in the laboratory:

\begin{itemize}
    \item \textbf{Learning from Experience}: If a particular catalyst performs well, you'll try similar catalysts (exploitation).
    \item \textbf{Exploring New Possibilities}: You'll also test some conditions that look different, in case you miss something better (exploration).
    \item \textbf{Iterative Refinement}: Each round of experiments informs the next, gradually converging toward optimal conditions.
\end{itemize}

\subsection{Mathematical Foundation}

Bayesian Optimization consists of two key components:

\subsubsection{Surrogate Model}

A probabilistic model that approximates the true objective function $f(x)$, where $x$ represents reaction conditions. Common choices include:

\begin{itemize}
    \item \textbf{Gaussian Process (GP)}: Provides uncertainty estimates alongside predictions
    \item \textbf{Random Forest (RF)}: Handles categorical variables effectively
    \item \textbf{Bayesian Neural Network (BNN)}: Captures complex nonlinear relationships
\end{itemize}

The surrogate model is updated iteratively as new experimental data becomes available.

\subsubsection{Acquisition Function}

A criterion that determines which experiment to perform next by balancing exploration and exploitation. Common acquisition functions include:

\begin{itemize}
    \item \textbf{EHVI} (Expected Hypervolume Improvement): Default for multi-objective optimization
    \item \textbf{UCB} (Upper Confidence Bound): Conservative strategy
    \item \textbf{ParEGO}: Transforms multi-objective into single-objective problems
    \item \textbf{NEI} (Noisy Expected Improvement): Accounts for experimental noise
\end{itemize}

% ============================================================================
\section{SynBO Architecture}
\label{sec:architecture}
% ============================================================================

\subsection{Project Structure}

SynBO is organized into modular components, each serving a specific purpose:

\begin{lstlisting}[language=bash, caption=SynBO Project Structure]
synbo/
├── src/synbo/
│   ├── synbo.py              # Main optimizer class
│   ├── initialize.py         # Initial sampling strategies
│   ├── optimize.py           # Optimization algorithm dispatcher
│   ├── cli.py                # Command-line interface
│   ├── algorithm/
│   │   ├── bo_core.py        # Bayesian optimization core
│   │   ├── acq_function.py   # Acquisition functions
│   │   ├── sg_model.py       # Surrogate models
│   │   ├── evolution.py      # Evolutionary algorithm
│   │   └── particle_swarm.py # Particle swarm algorithm
│   ├── descriptor/           # Molecular descriptor processing
│   ├── analysis/             # LLM-powered analysis module
│   └── utils/                # Utility functions
├── scripts/
│   ├── initialize.py         # Initialization script
│   ├── optimize.py           # Optimization script
│   └── get_desc.py          # Descriptor generation
└── examples/                 # Example workflows
\end{lstlisting}

\subsection{Core Components}

\subsubsection{ReactionOptimizer Class}

The \texttt{ReactionOptimizer} class is the heart of SynBO, providing a unified interface for:

\begin{itemize}
    \item Loading reaction space definitions
    \item Processing molecular descriptors
    \item Running initialization and optimization
    \item Managing experimental results
    \item Generating recommendations
\end{itemize}

\subsubsection{Descriptor Processing}

SynBO supports multiple descriptor types:

\begin{itemize}
    \item \textbf{RDKit Descriptors}: 200+ molecular fingerprints and physicochemical properties
    \item \textbf{DFT Calculated Descriptors}: Quantum mechanical properties
    \item \textbf{One-Hot Encoding}: For categorical variables without molecular structure
    \item \textbf{Custom Descriptors}: User-provided numerical representations
\end{itemize}

\subsubsection{LLM-Powered Analysis}

A unique feature of SynBO is its integration with Large Language Models (LLMs) to analyze failed experiments:

\begin{itemize}
    \item Identifies problematic reagent combinations
    \item Suggests chemical explanations for failures
    \item Automatically excludes ``problematic'' conditions from future recommendations
\end{itemize}

% ============================================================================
\section{Practical Workflow for Chemists}
\label{sec:workflow}
% ============================================================================

\subsection{Overview}

SynBO provides a skill-based interface that guides chemists through the optimization process in a structured manner. The workflow consists of four main phases:

\begin{enumerate}
    \item Project Setup and Configuration
    \item Reaction Space Definition
    \item Descriptor Generation
    \item Iterative Optimization
\end{enumerate}

\subsection{Phase 1: Project Setup}

Before starting optimization, users must configure the project environment:

\subsubsection{Working Directory and Project Name}

The system maintains a \texttt{config.json} file storing:
\begin{lstlisting}[language=json]
{
    "project_wd": "/path/to/project",
    "project_name": "my_reaction_optimization"
}
\end{lstlisting}

\textbf{Validation}: The system verifies that the working directory exists and sanitizes the project name (replacing spaces and special characters with underscores).

\subsubsection{Optimization Metrics Definition}

Users specify what they want to optimize through an \texttt{optimization\_settings.json} file:

\begin{lstlisting}[language=json]
{
    "optimization_settings": {
        "opt_metrics": ["yield", "ee"],
        "opt_direct_info": [
            {
                "opt_direct": "max",
                "opt_range": [0, 100],
                "metric_weight": 1.0
            },
            {
                "opt_direct": "max",
                "opt_range": [0, 100],
                "metric_weight": 2.0
            }
        ]
    }
}
\end{lstlisting}

This example optimizes both yield and enantioselectivity (ee), with ee given higher priority (weight = 2.0).

\subsection{Phase 2: Reaction Space Definition}

The reaction space defines all possible condition combinations. Users can provide this data via:

\subsubsection{Direct Input}

Provide SMILES strings for each reagent category directly in the chat interface.

\subsubsection{File Upload}

Upload CSV files with the following structure:

\begin{lstlisting}[language=csv]
SMILES,name
CCO,ethanol
CCCO,propanol
CCCOH,isopropanol
\end{lstlisting}

\textbf{Naming Convention}: If no names are provided, the system automatically assigns names using the format \texttt{\{reagent\_type\}-1}, \texttt{\{reagent\_type\}-2}, etc.

\textbf{Storage}: Each reagent type is saved as \texttt{\{reagent\_type\}.csv} in the \texttt{project\_wd/rxn\_space} directory.

\subsection{Phase 3: Descriptor Generation}

Condition descriptors are numerical representations of molecules used by the optimization algorithm.

\subsubsection{Automated RDKit Descriptor Generation}

If users request automated generation, the system:

\begin{enumerate}
    \item Verifies that reaction space data exists
    \item Confirms which reagent categories require descriptor calculation
    \item Invokes \texttt{scripts/get\_desc.py} to compute descriptors
\end{enumerate}

\subsubsection{Handling Invalid SMILES}

Entries with invalid SMILES (e.g., \texttt{no\_catalyst}, \texttt{none}, \texttt{blank}) are automatically handled:

\begin{itemize}
    \item Detected using RDKit's \texttt{MolFromSmiles} parser
    \item Skipped during descriptor calculation
    \item Added as placeholder rows with all descriptor values set to 0.0
\end{itemize}

This strategy ensures that categorical ``none'' options (e.g., ``no catalyst'', ``no additive'') are represented as zero vectors, allowing the optimizer to learn their effect independently.

\subsubsection{Continuous Variables}

For non-molecular variables (temperature, time, voltage), descriptors are formatted as:

\begin{lstlisting}[language=csv]
Name,value
temperature,25
temperature,50
temperature,80
time,12
time,24
\end{lstlisting}

\subsection{Phase 4: Initialization}

The initialization phase generates the first batch of recommended experiments when no previous data exists.

\subsubsection{Command}

\begin{lstlisting}[language=bash]
python scripts/initialize.py --desc-dir <path_to_descriptors> \
    --output <output_directory> --batch-size 8
\end{lstlisting}

\subsubsection{Key Parameters}

\begin{table}[h]
\centering
\caption{Initialization Parameters}
\begin{tabular}{lll}
\toprule
\textbf{Parameter} & \textbf{Default} & \textbf{Description} \\
\midrule
\texttt{--batch-size} & 5 & Number of initial samples \\
\texttt{--sampling-method} & kmeans & Sampling strategy \\
\texttt{--desc-normalize} & minmax & Descriptor normalization \\
\texttt{--refine-desc} & filter\_0.8 & Descriptor refinement \\
\texttt{--random-seed} & 42 & Reproducibility seed \\
\bottomrule
\end{tabular}
\end{table}

\subsubsection{Sampling Methods}

\begin{itemize}
    \item \textbf{Sobol}: Quasi-random low-discrepancy sequence
    \item \textbf{Random}: Pure random sampling
    \item \textbf{LHS} (Latin Hypercube): Stratified sampling
    \item \textbf{K-means}: Clustering-based sampling
\end{itemize}

\subsubsection{Output}

The system generates an Excel file (\texttt{recommended\_conditions.xlsx}) containing the first batch of recommended experimental conditions.

\subsection{Phase 5: Iterative Optimization}

After completing the initial experiments in the lab, users upload results and continue optimization.

\subsubsection{Input Data Format}

The results CSV file must contain:
\begin{itemize}
    \item Reagent type columns (matching reaction space)
    \item Optimization metric columns (e.g., \texttt{yield}, \texttt{ee})
    \item Optional: \texttt{batch} column indicating batch numbers
\end{itemize}

\begin{lstlisting}[language=csv]
catalyst,solvent,base,temperature,yield,ee
Pd(OAc)2,THF,Cs2CO3,25,45.2,82.1
Pd(PPh3)4,Dioxane,K2CO3,50,67.8,91.3
\end{lstlisting}

\subsubsection{Optimization Command}

\begin{lstlisting}[language=bash]
python scripts/optimize.py --desc-dir <path_to_descriptors> \
    --input experimental_results.csv \
    --output <output_directory> --batch-size 5
\end{lstlisting}

\subsubsection{Optimization Parameters}

\begin{table}[h]
\centering
\caption{Optimization Parameters}
\begin{tabular}{lll}
\toprule
\textbf{Parameter} & \textbf{Default} & \textbf{Description} \\
\midrule
\texttt{--batch-size} & 1 & Recommendations per round \\
\texttt{--optimize-method} & default\_BO & Optimization algorithm \\
\texttt{--surrogate-model} & GP & Surrogate model type \\
\texttt{--desc-normalize} & zscore & Descriptor normalization \\
\bottomrule
\end{tabular}
\end{table}

\subsubsection{Optimization Algorithms}

\begin{itemize}
    \item \textbf{default\_BO}: Standard Bayesian Optimization (recommended)
    \item \textbf{particle\_swarm}: Particle Swarm Optimization
    \item \textbf{evolution}: Evolutionary Algorithm
    \item \textbf{random\_select}: Random Search (baseline)
\end{itemize}

\subsubsection{Output Summary}

After optimization, the system displays:
\begin{itemize}
    \item Number of recommended new conditions
    \item Count of \textit{exploit} vs \textit{explore} recommendations
    \item Excel file with detailed conditions
\end{itemize}

\subsection{Phase 6: LLM-Powered Constraint Discovery}

When certain condition combinations repeatedly fail, SynBO can analyze patterns and suggest constraints:

\begin{lstlisting}[language=python]
# After round 3, let AI analyze which conditions to avoid
constraints = optimizer.get_constraints(method='llm')

# Apply constraints to next round of optimization
optimizer.run(batch_size=5, constraints=constraints)
\end{lstlisting}

\textbf{Application Scenarios}:
\begin{itemize}
    \item Discover ``DBU + high temperature'' always leads to decomposition $\rightarrow$ Auto-exclude
    \item Discover ``toluene solvent'' works best with specific catalyst $\rightarrow$ Prioritize similar combinations
\end{itemize}

\subsection{Phase 7: Progress Tracking}

For multi-objective optimization, SynBO provides Hypervolume metrics to track progress:

\begin{lstlisting}[language=python]
# Calculate Hypervolume for current Pareto front
hv = optimizer.calculate_current_hv()
print(f"Current optimization progress: {hv['hv_normalized']*100:.1f}%")

# View progress across rounds
progress = optimizer.calculate_hv_by_batch()
\end{lstlisting}

\textbf{Chemistry Interpretation}:
\begin{itemize}
    \item Hypervolume measures the ``performance space'' covered by currently found optimal conditions
    \item When Hypervolume growth slows down, you're near optimal and can consider stopping experiments
\end{itemize}

% ============================================================================
\section{Algorithms and Chemical Interpretation}
\label{sec:algorithms}
% ============================================================================

\subsection{Surrogate Models: Predicting Reaction Outcomes}

\begin{table}[h]
\centering
\caption{Surrogate Models and Their Chemical Intuition}
\begin{tabular}{p{3.5cm}p{5cm}p{4cm}}
\toprule
\textbf{Model} & \textbf{Chemistry Intuition} & \textbf{Best For} \\
\midrule
\textbf{Gaussian Process (GP)} & Assumes similar conditions give similar results & Fewer experiments (<50), clear reaction mechanisms \\
\textbf{Random Forest (RF)} & Voting via multiple decision trees & Many categorical variables (many catalyst/solvent types) \\
\textbf{BNN (Neural Network)} & Deep learning for complex nonlinear relationships & Large-scale screening (>100 experiments) \\
\textbf{Bayesian Linear} & Linear approximation, fast but simple & Preliminary screening, quick results \\
\bottomrule
\end{tabular}
\end{table}

\subsection{Acquisition Functions: Choosing the Next Experiment}

\begin{table}[h]
\centering
\caption{Acquisition Functions and Chemical Strategies}
\begin{tabular}{p{2.5cm}p{5.5cm}p{4.5cm}}
\toprule
\textbf{Function} & \textbf{Chemistry Strategy} & \textbf{When to Use} \\
\midrule
\textbf{EHVI} (Default) & Balance yield and selectivity, find Pareto optimal frontier & Optimizing yield and ee simultaneously \\
\textbf{UCB} & Conservative strategy, prioritize high-yield conditions with certainty & Limited time, cannot afford failures \\
\textbf{ParEGO} & Transform multi-objective into single-objective & More than 2 objectives (yield + ee + cost) \\
\textbf{NEI} & Account for experimental error & High variability in replicate experiments \\
\bottomrule
\end{tabular}
\end{table}

% ============================================================================
\section{Case Studies}
\label{sec:case_studies}
% ============================================================================

\subsection{Case 1: Asymmetric Hydrogenation}

\textbf{Background}: Screening chiral phosphoric acid catalysts for imine asymmetric hydrogenation.

\textbf{Reaction Space}:
\begin{itemize}
    \item 6 catalysts (CPA-1 through CPA-6)
    \item 4 additives (MsOH, TfOH, TFA, None)
    \item 4 solvents (DCE, PhCF3, Toluene, Et2O)
    \item 4 temperatures (-20, 0, 25, 40 $^\circ$C)
    \item 4 H$_2$ pressures (1, 10, 20, 50 atm)
\end{itemize}

\textbf{Total Combinations}: $6 \times 4 \times 4 \times 4 \times 4 = 1536$

\textbf{Optimization Objective}: High yield + High ee

\textbf{Result}: Only 24 experiments needed (vs. 1536 full combinations), found conditions with 94\% yield and 98\% ee.

\subsection{Case 2: Buchwald-Hartwig Amination}

\textbf{Background}: Pd-catalyzed aromatic amination, screening ligand and base combinations.

\textbf{LLM Analysis}: After several rounds, SynBO's LLM module identified that ``XPhos + strong base'' combinations consistently led to catalyst deactivation.

\textbf{Outcome}: These combinations were automatically excluded from future recommendations, saving significant experimental time and resources.

% ============================================================================
\section{Best Practices and Tips}
\label{sec:best_practices}
% ============================================================================

\subsection{Experimental Design}

\begin{enumerate}
    \item \textbf{Initial Batch Size}: Start with 5--10 experiments for initialization
    \item \textbf{Subsequent Batches}: Use 1--5 experiments per optimization round
    \item \textbf{Replicates}: Include 2--3 replicate experiments to estimate experimental noise
    \item \textbf{Diversity}: Ensure initial conditions cover the full reaction space
\end{enumerate}

\subsection{Descriptor Selection}

\begin{itemize}
    \item \textbf{RDKit Descriptors}: Suitable for most organic molecules
    \item \textbf{DFT Descriptors}: Use when electronic properties are critical
    \item \textbf{One-Hot Encoding}: For simple categorical variables
    \item \textbf{Hybrid Approach}: Combine multiple descriptor types for complex systems
\end{itemize}

\subsection{When to Stop}

Monitor the following indicators:
\begin{itemize}
    \item \textbf{Hypervolume Plateau}: When HV growth slows significantly
    \item \textbf{Yield Saturation}: When recommended conditions consistently give similar yields
    \item \textbf{Resource Constraints}: When time or material budget is exhausted
    \item \textbf{Target Achievement}: When desired yield/selectivity is reached
\end{itemize}

% ============================================================================
\section{Conclusion and Future Perspectives}
\label{sec:conclusion}
% ============================================================================

\subsection{Summary}

SynBO represents a paradigm shift in how synthetic chemists approach reaction optimization. By combining Bayesian Optimization with an intuitive skill-based interface, it enables chemists to:

\begin{itemize}
    \item Reduce experimental effort by 80--90\%
    \item Systematically explore complex reaction spaces
    \item Learn from failures through LLM-powered analysis
    \item Track optimization progress quantitatively
\end{itemize}

\subsection{Future Directions}

Potential enhancements include:
\begin{itemize}
    \item Integration with automated laboratory platforms
    \item Support for more complex reaction types (cascade, tandem)
    \item Enhanced descriptor libraries (3D structures, dynamics)
    \item Community-driven model sharing and transfer learning
\end{itemize}

\subsection{Availability}

SynBO is available as an open-source Python package:
\begin{itemize}
    \item \textbf{Installation}: \texttt{pip install synbo}
    \item \textbf{Documentation}: Available in the \texttt{docs/} directory
    \item \textbf{Examples}: Comprehensive examples in \texttt{examples/} directory
    \item \textbf{License}: MIT License
\end{itemize}

% ============================================================================
\section*{Acknowledgments}
% ============================================================================

The author thanks the synthetic chemistry community for valuable feedback and suggestions that shaped SynBO's development.

% ============================================================================
\begin{thebibliography}{99}
% ============================================================================

\bibitem{shahriari2016}
Shahriari, B., Swersky, K., Wang, Z., Adams, R. P., \& de Freitas, N. (2016).
Taking the human out of the loop: A review of Bayesian optimization.
\textit{Proceedings of the IEEE}, 104(1), 148--175.

\bibitem{frazier2018}
Frazier, P. I. (2018).
A tutorial on Bayesian optimization.
\textit{arXiv preprint arXiv:1807.02811}.

\bibitem{hernandez2017}
Hernández-Lobato, J. M., Requeima, J., Pyzer-Knapp, E. O., \& Aspuru-Guzik, A. (2017).
Parallel and distributed Thompson sampling for large-scale accelerated exploration of chemical space.
\textit{International Conference on Machine Learning}, 1470--1479.

\bibitem{gregoire2019}
Grégoire, J. M., \& Aspuru-Guzik, A. (2019).
Accelerated materials discovery with synthetic intelligence.
\textit{ACS Central Science}, 5(4), 601--603.

\bibitem{shields2021}
Shields, B. J., Stevens, J., Li, J., et al. (2021).
Bayesian reaction optimization as a tool for chemical synthesis.
\textit{Nature}, 590(7844), 89--96.

\bibitem{tan2025}
Tan, Z. (2025).
SynBO: Synthetic Bayesian Optimization for Chemical Reaction Optimization.
\textit{GitHub Repository}. \url{https://github.com/yourusername/synbo}

\end{thebibliography}

% ============================================================================
\appendix
% ============================================================================

\section{Installation Guide}

\subsection{System Requirements}

\begin{itemize}
    \item Python 3.12 or higher
    \item 4 GB RAM minimum (8 GB recommended)
    \item 1 GB free disk space
\end{itemize}

\subsection{Installation Steps}

\begin{lstlisting}[language=bash]
# Create virtual environment (recommended)
python -m venv synbo_env
source synbo_env/bin/activate  # On Windows: synbo_env\Scripts\activate

# Install SynBO
pip install synbo

# Verify installation
python -c "from synbo import ReactionOptimizer; print('SynBO installed successfully!')"
\end{lstlisting}

\subsection{Optional Dependencies}

\begin{lstlisting}[language=bash]
# For RDKit descriptor generation
pip install rdkit

# For LLM-powered analysis
pip install openai

# For visualization
pip install matplotlib seaborn
\end{lstlisting}

\section{Quick Reference Card}

\subsection{Common Commands}

\begin{lstlisting}[language=bash]
# Initialize optimization (no previous data)
python scripts/initialize.py --desc-dir ./descriptors --batch-size 8

# Continue optimization (with experimental results)
python scripts/optimize.py --desc-dir ./descriptors --input results.csv --batch-size 5

# Generate RDKit descriptors
python scripts/get_desc.py --input ./rxn_space --output ./descriptors
\end{lstlisting}

\subsection{File Structure}

\begin{lstlisting}[language=bash]
project_directory/
├── config.json                    # Project configuration
├── optimization_settings.json     # Optimization metrics
├── rxn_space/                     # Reaction space definitions
│   ├── catalyst.csv
│   ├── solvent.csv
│   └── base.csv
├── mol_desc/                      # Molecular descriptors
│   ├── catalyst_RDKit.csv
│   ├── solvent_RDKit.csv
│   └── base_RDKit.csv
└── output/                        # Optimization results
    ├── batch_0/
    └── batch_1/
\end{lstlisting}

\end{document}